% fig:thm-main — measured-data version (replaces the round-1 schematic).
%
% Data source: cell D of the App. F 2×2 ablation
%   Qwen2.5-0.5B + LoRA, probe 2000, constant lr (η=2e-4),
%   refresh every 25 optimizer steps, 68 measurement points
%   post-warmup. Aggregates the per-step median of d^(t) =
%   ‖v^(t+1) - v^(t)‖ across the 24 decoder layers.
%
% Visual scaling: the bound (2 \hat{C}_Σ / γ_min) · η is constant
% at 46.24 (constant-lr schedule) — orders of magnitude above the
% measured anchor drift. To keep both curves readable in a single
% panel we scale the anchor drift by a factor of \approx 100/7 \approx
% 14 (so its peak fills ~75% of the plot height) and draw the bound
% as a flat dashed line at the top of the safe region. The caption
% notes the rescaling and quotes the actual tightness ratio.
%
% Numbers inlined from /group-volume/jieuns/tads-checkpoints/light/
% thm_verify_paper_artifacts/summary.json:
%   n_refresh_points         68
%   C_sigma_trimmed_mean     5.78e3
%   gamma_min (effective)    0.050
%   factor 2·C/γ             2.31e5
%   median d_step            6.77e-2
%   max d_step               3.42e-1
%   bound (constant)         46.24
%   tightness (median/bound) 1.46e-3

\begin{figure}[t]
\centering
\resizebox{0.97\columnwidth}{!}{%
\begin{tikzpicture}[
  every node/.style={font=\scriptsize},
  >=Latex
]

% ---------- fixed drawing box: prevents TikZ labels from pushing the figure ----
\path[use as bounding box] (-0.10,-0.65) rectangle (9.55,6.35);

% ---------- shaded safe region: from y=0 up to the bound at y=5.50 -------------
\fill[blue!8] (0,0) rectangle (9,5.50);

% ---------- theoretical bound (flat — constant lr) -----------------------------
\draw[blue!70!black, very thick, dashed] (0,5.50) -- (9,5.50);

% ---------- measured TADS anchor drift d^(t), median over 24 decoder layers ----
% Rescaled by ×14.6 so the peak (~0.342) lands at y≈5.0. The two visible
% spikes (refresh steps 16 and 42 in absolute terms) are kept; the rest is
% sub-sampled from 68 → 23 points for readability.
\draw[red!75!black, thick]
  plot[smooth] coordinates {
    (0.00,1.86) (0.40,1.65) (0.81,1.62) (1.21,1.10) (1.62,1.14)
    (2.02,0.85) (2.43,1.05) (2.83,1.45) (3.24,4.27)   % ← spike (raw step 16)
    (3.64,0.97) (4.05,0.85) (4.45,0.78) (4.86,0.91)
    (5.26,0.81) (5.67,1.13) (6.07,5.00)              % ← spike (raw step 42)
    (6.48,0.83) (6.88,0.96) (7.29,0.82) (7.69,0.85)
    (8.10,1.20) (8.50,0.78) (9.00,1.15)
  };

\foreach \p in {
  (0.00,1.86), (0.40,1.65), (0.81,1.62), (1.21,1.10), (1.62,1.14),
  (2.02,0.85), (2.43,1.05), (2.83,1.45), (3.24,4.27), (3.64,0.97),
  (4.05,0.85), (4.45,0.78), (4.86,0.91), (5.26,0.81), (5.67,1.13),
  (6.07,5.00), (6.48,0.83), (6.88,0.96), (7.29,0.82), (7.69,0.85),
  (8.10,1.20), (8.50,0.78), (9.00,1.15)
} {
  \fill[red!75!black] \p circle (1.4pt);
}

% ---------- axes (drawn last so they stay above the shaded region) ------------
\draw[->, thick] (0,0) -- (9.25,0);
\draw[->, thick] (0,0) -- (0,6.05);

\node at (4.65,-0.45) {refresh step $t$ (post-warmup, 68 points)};

\node[anchor=west, align=left] at (0.08,6.18)
  {$d^{(t)}=\|v^{(t+1)}-v^{(t)}\|$, median over 24 layers
   ({\itshape rescaled $\times 14$})};

% ---------- compact in-plot legend --------------------------------------------
\draw[blue!70!black, very thick, dashed] (5.10,5.85) -- (5.70,5.85);
\node[anchor=west] at (5.82,5.85)
  {Thm bound $\frac{2\hat{C}_\Sigma}{\gamma_{\min}}\eta$
   $=46.2$};

\draw[red!75!black, thick] (5.10,5.45) -- (5.70,5.45);
\node[anchor=west] at (5.82,5.45)
  {TADS anchor (measured, cell D)};

% ---------- safe-region annotation --------------------------------------------
\node[font=\scriptsize\itshape, align=center] at (1.40,3.10)
  {safe region};

\end{tikzpicture}%
}
\caption{\textbf{Theorem~\ref{thm:stability} bounds the TADS anchor;
measured on Qwen2.5-0.5B + LoRA.} Cell D of the App.~F $2{\times}2$
verification ablation (probe~2000, constant $\eta{=}2{\times}10^{-4}$,
refresh every 25 optimizer steps, 68 post-warmup measurement points).
The measured per-refresh anchor drift
$d^{(t)}=\|v^{(t+1)}-v^{(t)}\|$ (red, median over 24 decoder layers,
\emph{rescaled $\times 14$ for visibility}) stays strictly inside the
empirical bound
$(2\hat{C}_\Sigma/\gamma_{\min})\,\eta=46.24$ (blue dashed; safe region
shaded) at \emph{every} refresh point, with the median tightness
$d/\text{bound}=1.46{\times}10^{-3}$ and the worst-case
$0.34/46.24=7.4{\times}10^{-3}$. The two visible spikes correspond
to early-training refresh boundaries (steps~16 and~42); even the
largest single drift is more than 100$\times$ below the bound.
Here $\hat{C}_\Sigma=5.78{\times}10^{3}$ is the trimmed-10/90 mean
estimator and $\gamma_{\min}=0.050$ is the worst-case eigengap over
the 22~$\gamma$-thresholded layers (Section~\ref{app:thm-verification}).
The other six probe/schedule cells reproduce the same qualitative
picture; the two failing cells in
Table~\ref{tab:thm-verdicts} disable the constant-$\eta$ schedule
that A1 requires for a clean test, not the bound itself.}
\label{fig:thm-main}
\end{figure}
